% from ega2-1-fr.tex
\oldpage[II]{77}

de $Q$, où $B=R_{(f)}\otimes_AS_{(g)}$~; la restriction de $\mathcal{L}$ à
cet ouvert affine est $\widetilde{L}$, où
\[
  L=(R(1))_{(f)}\otimes_A(S(1))_{(g)}
\]
et l'élément $c=(f/1)\otimes(g/1)$ est un générateur de $L$ considéré comme
$B$-module libre (\hyperref[II.2.5.7-fr]{2.5.7}). L'homomorphisme
(\hyperref[II.4.3.1.1-fr]{4.3.1.1}) correspond à l'homomorphisme
\[
  (x\otimes y)\otimes b\longmapsto b((x/1)\otimes(y/1))
\]
de $(E\otimes_A F)\otimes_AB$ dans $L$. Avec les notations de
(\hyperref[II.3.7.2-fr]{3.7.2}), on a donc
$v_1(x\otimes y)=(x/f)\otimes(y/g)$~; la restriction de $\varsigma$ à
$D_+(f)\times_YD_+(g)$ est un morphisme de ce schéma affine dans
$D_+(f\otimes g)$, correspondant à l'homomorphisme d'anneaux
$\omega:T_{(f\otimes g)}\to R_{(f)}\otimes_AS_{(g)}$ tel que
\[
\label{II.4.3.1.3-fr}
  \omega((x\otimes y)/(f\otimes g))=(x/f)\otimes(y/g)
\tag{4.3.1.3}
\]
pour $x\in E$ et $y\in F$.
\end{env}

\begin{env}[4.3.2]
\label{II.4.3.2-fr}
Il résulte de (\hyperref[II.4.2.3-fr]{4.2.3}) que l'on a un isomorphisme
canonique
\[
\label{II.4.3.2.1-fr}
  \tau:\varsigma^*(\mathcal{O}_P(1))\xrightarrow{\sim}
  \mathcal{O}_{P_1}(1)\otimes_Y\mathcal{O}_{P_2}(1)
\tag{4.3.2.1}
\]
où on a posé
$P=\mathbf{P}(\mathcal{E}\otimes_{\mathcal{O}_Y}\mathcal{F})$. Montrons
que, pour $x\in\Gamma(Y,\mathcal{E})$ et
$y\in\Gamma(Y,\mathcal{F})$, on a
\[
\label{II.4.3.2.2-fr}
  \tau(\alpha_1(x\otimes y))=\alpha_1(x)\otimes\alpha_1(y).
\tag{4.3.2.2}
\]
En effet, on est ramené au cas où $Y$ est affine, et on a alors, avec les
notations de (\hyperref[II.4.3.1-fr]{4.3.1}) et de
(\hyperref[II.2.6.2-fr]{2.6.2}),
$\alpha_1^{f\otimes g}(x\otimes y)=(x\otimes y)/1$ dans
$(T(1))_{(f\otimes g)}$, $\alpha_1^f(x)=x/1$ dans $(R(1))_{(f)}$ et
$\alpha_1^g(y)=y/1$ dans $(S(1))_{(g)}$. La définition de $\tau$ donnée dans
(\hyperref[II.4.2.3-fr]{4.2.3}) et le calcul de $v_1$ fait dans
(\hyperref[II.4.3.1-fr]{4.3.1}) prouvent alors aussitôt
(\hyperref[II.4.3.2.2-fr]{4.3.2.2}). On en déduit la formule
\[
\label{II.4.3.2.3-fr}
  \varsigma^{-1}(P_{x\otimes y})=(P_1)_x\times_Y(P_2)_y
\tag{4.3.2.3}
\]
avec les notations de (\hyperref[II.3.1.4-fr]{3.1.4}). En effet, compte tenu
de (\hyperref[II.3.3.3-fr]{3.3.3}), la formule
(\hyperref[II.4.3.2.2-fr]{4.3.2.2}) ramène (en revenant au cas affine, à
l'aide de (\hyperref[I.3.2.7-fr]{I, 3.2.7} et
(\hyperref[I.3.2.3-fr]{3.2.3})) à prouver le lemme suivant~:

\begin{lemma}[4.3.2.4]
\label{II.4.3.2.4-fr}
Soient $B$, $B'$ deux $A$-algèbres, et soient
$Y=\operatorname{Spec}(A)$, $Z=\operatorname{Spec}(B)$,
$Z'=\operatorname{Spec}(B')$~; quels que soient $t\in B$, $t'\in B'$, on a
$D(t\otimes t')=D(t)\times_YD(t')$.
\end{lemma}

\begin{proof}
En effet, si $p:Z\times_YZ'\to Z$ et $p':Z\times_YZ'\to Z'$ sont les
projections canoniques, il résulte de
(\hyperref[I.1.2.2.2-fr]{I, 1.2.2.2}) que l'on a
$p^{-1}(D(t))=D(t\otimes1)$ et
$p'^{-1}(D(t'))=D(1\otimes t')$~; la conclusion résulte de
(\hyperref[I.3.2.7-fr]{I, 3.2.7}) et
(\hyperref[I.1.1.9.1-fr]{I, 1.1.9.1}), puisque
$(t\otimes1)(1\otimes t')=t\otimes t'$.
\end{proof}
\end{env}

\begin{proposition}[4.3.3]
\label{II.4.3.3-fr}
Le morphisme de Segre est une immersion fermée.
\end{proposition}

\begin{proof}
En effet, la question étant locale sur $Y$, on est ramené au cas où $Y$ est
affine. Avec les notations de (\hyperref[II.4.3.1-fr]{4.3.1}) et
(\hyperref[II.4.3.2-fr]{4.3.2}), les $D_+(f\otimes g)$ forment alors une base
de la topologie de $P$, puisque les $f\otimes g$ engendrent $T$ lorsque $f$
parcourt $E$ et $g$ parcourt $F$. D'autre part, on a en vertu de
(\hyperref[II.4.3.2.3-fr]{4.3.2.3})
$\varsigma^{-1}(D_+(f\otimes g))=D_+(f)\times_YD_+(g)$. Il suffit donc
(\hyperref[I.4.2.4-fr]{I, 4.2.4}) de prouver que la restriction de
$\varsigma$ à $D_+(f)\times_YD_+(g)$ est une immersion fermée dans
$D_+(f\otimes g)$. Or, avec les mêmes notations, la formule
(\hyperref[II.4.3.1.3-fr]{4.3.1.3}) montre que $\omega$ est surjectif, ce qui
termine la démonstration.
\end{proof}

\begin{env}[4.3.4]
\label{II.4.3.4-fr}
Le morphisme de Segre est fonctoriel en $\mathcal{E}$ et $\mathcal{F}$,
lorsqu'on restreint les

